% 2013 湖南卷文科第 17 题：直棱柱 ABC-A1B1C1，D 为 BC 中点，E 在 BB1 上。
\begin{tikzpicture}[font=\footnotesize, >=Stealth, line width=0.45pt, scale=1.05]
  \coordinate (B)  at (0,0);
  \coordinate (C)  at (0.5,2.6);
  \coordinate (C1) at (4.1,2.6);
  \coordinate (B1) at (3.6,0);
  \coordinate (D)  at (0.25,1.3);
  \coordinate (A)  at (1.6,0.92);
  \coordinate (A1) at (5.2,0.92);
  \coordinate (E)  at (1.51,0);
  % 前面 BCC1B1
  \draw (B) -- (C) -- (C1) -- (B1) -- cycle;
  % 上底边与 C1E
  \draw (C1) -- (A1);
  \draw (C1) -- (E);
  % 被遮挡的棱
  \draw[dashed] (B) -- (A);
  \draw[dashed] (C) -- (A);
  \draw[dashed] (A) -- (A1);
  \draw[dashed] (A1) -- (B1);
  % 点
  \fill (A) circle (1.4pt) (A1) circle (1.4pt) (B) circle (1.4pt)
        (B1) circle (1.4pt) (C) circle (1.4pt) (C1) circle (1.4pt)
        (D) circle (1.4pt) (E) circle (1.4pt);
  % 标注
  \node[above left]  at (C)  {$C$};
  \node[above right] at (C1) {$C_1$};
  \node[below right] at (A)  {$A$};
  \node[right]       at (A1) {$A_1$};
  \node[below left]  at (B)  {$B$};
  \node[below right] at (B1) {$B_1$};
  \node[left]        at (D)  {$D$};
  \node[below]       at (E)  {$E$};
\end{tikzpicture}
